%============================================================
% CH1 · P3 · PROBLEM 1 — ROTATE ARRAY (LeetCode 189)
%============================================================
\subsection{Problem 1.3.1: Rotate Array}
\problemlabel{1.3.1}{ssec:ch1-rotate-array}

\problemcard{1.3.1}{Rotate Array}{LeetCode 189}{https://leetcode.com/problems/rotate-array/}{\exmedium\ \lc{Medium}}{Rotation by three reversals ($AB \to BA$); the cycle method with $\gcd(n,k)$ cycles}

\subsubsection*{1. Problem statement}

You are given an integer array \texttt{nums} of length $n$ and a non-negative integer $k$. Rotate the array to the \textbf{right} by $k$ steps: in one step every element moves one position to the right and the last element moves to the front. The array must be modified \textbf{in place} (the function returns nothing). The follow-up asks for $\bigO{1}$ extra memory.

\textbf{Constraints:} $1 \le n \le 10^5$, \; $-2^{31} \le \texttt{nums[i]} \le 2^{31}-1$, \; $0 \le k \le 10^5$.

\begin{table}[H]
\centering\small
\begin{tabular}{@{}llcl@{}}
\toprule
\textbf{nums} & $k$ & \textbf{Result} & \textbf{Why} \\
\midrule
$[1,2,3,4,5,6,7]$ & 3 & $[5,6,7,1,2,3,4]$ & the last 3 values move to the front \\
$[-1,-100,3,99]$ & 2 & $[3,99,-1,-100]$ & the two halves swap places \\
\bottomrule
\end{tabular}
\end{table}

\begin{diagrambox}
\begin{tikzpicture}[x=1cm, y=1cm]
    \dpindex{i}{2.45}{index}{0}{6}
    \dparray{a}{1.8}{before}{1/acell, 2/acell, 3/acell, 4/acell, 5/aread, 6/aread, 7/aread}
    \dparray{b}{-0.2}{after}{5/aread, 6/aread, 7/aread, 1/akeep, 2/akeep, 3/akeep, 4/akeep}
    \foreach \x/\y in {0/3, 1/4, 2/5, 3/6}{ \draw[faded, successgreen!70!black] (a\x.south) -- (b\y.north); }
    \foreach \x/\y in {4/0, 5/1, 6/2}{ \draw[dep] (a\x.south) -- (b\y.north); }
    \node[note, anchor=west, text width=4.1cm, align=left] at (6.9,0.8)
        {index $i$ moves to $(i + 3) \bmod 7$: green values shift right by 3, orange values wrap around to the front};
\end{tikzpicture}
\figcap{Sample 1. Every element moves to $(i+k) \bmod n$; the order inside the green block and inside the orange block is unchanged.}
\end{diagrambox}

\subsubsection*{2. The problem model}
\begin{modelbox}[Problem Model --- Rotate Array]
\renewcommand{\arraystretch}{1.35}
\begin{tabularx}{\linewidth}{@{}>{\bfseries}l X@{}}
Input & an array $a[0..n-1]$ of 32-bit integers and a shift $k \ge 0$ \\
Computed & a new arrangement of the same array: $a'[(i + k) \bmod n] = a[i]$ for every $i$, written back into $a$ \\
Objects & positions of the array; the values are only moved, never compared or added (so no overflow is possible) \\
Size & $n \le 10^5$, $k \le 10^5$: $n \cdot k$ can reach $10^{10}$, too much; $\bigO{n}$ or $\bigO{n \log n}$ is needed \\
Memory & the follow-up: $\bigO{1}$ extra cells besides the array itself \\
\end{tabularx}
\end{modelbox}

\begin{diagrambox}
\begin{tikzpicture}[x=1cm, y=1cm,
    blkA/.style={rectangle, draw=primaryblue, fill=primaryblue!15, minimum height=7mm, font=\small\bfseries},
    blkB/.style={rectangle, draw=accentorange, fill=answerfill, minimum height=7mm, font=\small\bfseries}]
    \node[font=\small\bfseries, text=deepblue, anchor=west] at (-0.2,1.15) {the array as two blocks};
    \node[blkA, minimum width=3.6cm, anchor=west] (A) at (0,0.4) {$A = a[0..n-k-1]$};
    \node[blkB, minimum width=2.7cm, anchor=west] (B) at (3.6,0.4) {$B = $ last $k$};
    \draw[flow] (6.6,0.4) -- node[lbl, above]{rotate right by $k$} (8.6,0.4);
    \node[blkB, minimum width=2.7cm, anchor=west] at (8.8,0.4) {$B$};
    \node[blkA, minimum width=3.6cm, anchor=west] at (11.5,0.4) {$A$};
    \node[note] at (7.6,-0.55) {$|A| = n - k$, $|B| = k$ (after $k \leftarrow k \bmod n$)};
\end{tikzpicture}
\figcap{The whole problem in block language: $AB \to BA$. Nothing inside a block changes order.}
\end{diagrambox}

\subsubsection*{3. Key observations}

\begin{itemize}[leftmargin=*]
    \item \textbf{Only $k \bmod n$ matters.} $n$ steps bring every element back home. With $n = 1$ and $k = 10^5$ the answer is the input itself; code that uses $k$ directly as an index crashes.
    \item \textbf{The obvious in-place loop is wrong.} \texttt{for i: a[(i+k)\%n] = a[i]} overwrites $a[3]$ with $a[0]$ before $a[3]$ has been moved. On Sample 1 it produces $[2,3,1,1,2,3,1]$: values are lost and others are duplicated. Each value must be read \emph{before} its cell is overwritten.
    \item \textbf{A second array fixes that, but costs $\bigO{n}$ memory:} \texttt{b[(i+k)\%n] = a[i]}, then copy back. Correct and fast, but not the follow-up.
    \item \textbf{The block view gives an in-place method:} $AB \to BA$ is exactly what three reversals do (Tool 2 of the theory). The positions also form $\gcd(n,k)$ cycles $i \to (i+k) \bmod n$, which gives a second in-place method (Tool 3).
\end{itemize}

\begin{diagrambox}
\begin{tikzpicture}[x=1cm, y=1cm]
    \dparray{w}{1.1}{wrong loop}{2/aout, 3/aout, 1/aout, 1/aout, 2/aout, 3/aout, 1/aout}
    \dparray{r}{-0.1}{correct}{5/akeep, 6/akeep, 7/akeep, 1/akeep, 2/akeep, 3/akeep, 4/akeep}
    \node[note, anchor=west, text width=4.3cm, align=left] at (6.9,0.5)
        {\texttt{a[(i+3)\%7] = a[i]} for $i = 0..6$ on $[1..7]$: by the time $i = 3$ the cell already holds the old $a[0]$, and the values $4,5,6,7$ are lost for ever};
\end{tikzpicture}
\figcap{The counter-example to ``just move every value to its new place''. Writing in place needs an order of moves that never destroys an unread value.}
\end{diagrambox}

\subsubsection*{4. Approaches}

\begin{bruteforcebox}
\textbf{Idea.} Rotate by one step, $k$ times: remember the last element, shift everything one cell to the right (from the right end, so that nothing is overwritten too early), put the remembered element at index 0.
\begin{lstlisting}[style=cpp]
#include <bits/stdc++.h>
using namespace std;

class Solution {
public:
    void rotate(vector<int>& nums, int k) {
        int n = nums.size();
        for (int step = 0; step < k; step++) {     // k times: rotate by one
            int last = nums[n - 1];
            for (int i = n - 1; i > 0; i--) nums[i] = nums[i - 1];
            nums[0] = last;
        }
    }
};

int main() {                                       // input: n k, then the n values
    int n, k;
    if (!(cin >> n >> k)) return 0;
    vector<int> a(n);
    for (auto& x : a) cin >> x;
    Solution().rotate(a, k);
    for (int i = 0; i < n; i++) cout << a[i] << (i + 1 < n ? ' ' : '\n');
}
\end{lstlisting}
\begin{complexitybox}
\textbf{Time:} $\bigO{n \cdot k}$ --- up to $10^5 \cdot 10^5 = 10^{10}$ moves: \tle. (Using $k \bmod n$ helps only when $k \ge n$.) \textbf{Space:} $\bigO{1}$. It is the reference answer for the stress test below.
\end{complexitybox}
\end{bruteforcebox}

\begin{algorithmbox}
\textbf{Idea.} Reduce $k \leftarrow k \bmod n$. Reverse the whole array ($AB \to B^R A^R$), then reverse the first $k$ cells ($B^R \to B$) and the last $n-k$ cells ($A^R \to A$). Each reversal is the two-index loop of Pattern~3.
\begin{lstlisting}[style=cpp]
#include <bits/stdc++.h>
using namespace std;

class Solution {
    // reverse the half-open range [l, r) with two indices moving towards each other
    void rev(vector<int>& a, int l, int r) {
        for (r--; l < r; l++, r--) swap(a[l], a[r]);
    }
public:
    void rotate(vector<int>& nums, int k) {
        int n = nums.size();
        k %= n;                  // rotating by n changes nothing
        rev(nums, 0, n);         // 1. reverse everything
        rev(nums, 0, k);         // 2. reverse the first k (the old tail)
        rev(nums, k, n);         // 3. reverse the last n-k (the old head)
    }
};

int main() {                     // input: n k, then the n values
    int n, k;
    if (!(cin >> n >> k)) return 0;
    vector<int> a(n);
    for (auto& x : a) cin >> x;
    Solution().rotate(a, k);
    for (int i = 0; i < n; i++) cout << a[i] << (i + 1 < n ? ' ' : '\n');
}
\end{lstlisting}
\begin{complexitybox}
\textbf{Time:} $\bigO{n}$ --- the three reversals perform $\lfloor n/2\rfloor + \lfloor k/2\rfloor + \lfloor (n-k)/2\rfloor \le n$ swaps (Sample 1: $3 + 1 + 2 = 6$ swaps). \textbf{Space:} $\bigO{1}$ extra. \textbf{Verdict:} \ac{}.
\end{complexitybox}
\end{algorithmbox}

\begin{invariantbox}[Why three reversals rotate the array]
Let $m = n - k$ (after reduction, $0 \le k < n$). Take the element at old index $i$.
\begin{itemize}[leftmargin=*]
    \item If $i < m$ (block $A$): reverse all sends it to $n-1-i$, which is $\ge k$; reverse $[k, n)$ sends position $p$ to $k + (n-1) - p$, so it ends at $k + (n-1) - (n-1-i) = i + k$.
    \item If $i \ge m$ (block $B$): reverse all sends it to $n-1-i < k$; reverse $[0,k)$ sends $p$ to $k-1-p$, so it ends at $k-1-(n-1-i) = i + k - n = (i+k) \bmod n$.
\end{itemize}
Every element ends at $(i + k) \bmod n$, which is the definition of the rotation. Each reversal is correct by the invariant of Tool~1, and the second and third reversals touch disjoint ranges. \hfill$\blacksquare$
\end{invariantbox}

\begin{alternativebox}
\textbf{The cycle (juggling) method.} The move $i \to (i+k) \bmod n$ splits the positions into $g = \gcd(n,k)$ cycles of length $n/g$; the cycle of $s$ consists of the positions $\equiv s \pmod g$, so the starts $0, \dots, g-1$ reach every cycle once. Walk each cycle with one spare variable \texttt{carry}: drop the carried value at the next position and pick up the value that was there.
\begin{lstlisting}[style=cpp]
#include <bits/stdc++.h>
using namespace std;

class Solution {
public:
    void rotate(vector<int>& nums, int k) {
        int n = nums.size();
        k %= n;
        if (k == 0) return;
        int cycles = __gcd(n, k);                 // number of disjoint cycles
        for (int start = 0; start < cycles; start++) {
            int cur = start, carry = nums[start]; // the value travelling along the cycle
            do {
                int nxt = (cur + k) % n;          // where the value at cur must go
                swap(carry, nums[nxt]);           // drop it there, pick up the old value
                cur = nxt;
            } while (cur != start);
        }
    }
};

int main() {                                      // input: n k, then the n values
    int n, k;
    if (!(cin >> n >> k)) return 0;
    vector<int> a(n);
    for (auto& x : a) cin >> x;
    Solution().rotate(a, k);
    for (int i = 0; i < n; i++) cout << a[i] << (i + 1 < n ? ' ' : '\n');
}
\end{lstlisting}
Exactly $n$ writes, $\bigO{1}$ extra memory, $\bigO{n}$ time (plus $\bigO{\log n}$ for the gcd). Without the gcd one can instead count the moved elements and start a new cycle at $s+1$ until $n$ elements have moved --- equivalent, because the cycles are exactly the residue classes mod $g$.

\begin{diagrambox}
\begin{tikzpicture}[x=1cm, y=1cm]
    \dparray{c}{0}{positions}{0/gvA, 1/gvB, 2/gvA, 3/gvB, 4/gvA, 5/gvB}
    \foreach \x/\y in {0/2, 2/4}{ \draw[gde, draw=primaryblue, line width=1.1pt] (c\x.north) to[bend left=40] (c\y.north); }
    \draw[gde, draw=primaryblue, line width=1.1pt] (c4.north) to[bend right=25] (c0.north);
    \foreach \x/\y in {1/3, 3/5}{ \draw[gde, draw=recursionred, line width=1.1pt] (c\x.south) to[bend right=40] (c\y.south); }
    \draw[gde, draw=recursionred, line width=1.1pt] (c5.south) to[bend left=25] (c1.south);
    \node[note, anchor=west, text width=4.6cm, align=left] at (6.0,0)
        {$n = 6$, $k = 2$, $\gcd = 2$: blue cycle $0 \to 2 \to 4 \to 0$ (even positions), red cycle $1 \to 3 \to 5 \to 1$ (odd positions)};
\end{tikzpicture}
\figcap{Two starts are needed: one cycle never visits the positions of the other.}
\end{diagrambox}
\end{alternativebox}

\begin{tricksbox}
\begin{itemize}[leftmargin=*]
    \item \textbf{Library one-liner:} \texttt{std::rotate(nums.begin(), nums.end() - k, nums.end())} (after $k \bmod n$) makes the element at \texttt{end() - k} the new first element --- a right rotation by $k$, in place and linear.
    \item \textbf{Extra-array version} when memory does not matter: \texttt{vector<int> b(n); for i: b[(i + k) \% n] = nums[i]; nums = b;} --- the simplest code to get right in a contest.
    \item \textbf{Equivalent orders:} ``reverse $[0, n-k)$, reverse $[n-k, n)$, reverse all'' is also a right rotation by $k$ (it is the left rotation by $n-k$). Pick one form and remember it with a 3-element example.
    \item \textbf{Values are only moved}: no \texttt{long long} is needed even though they span the full \texttt{int} range.
\end{itemize}
\end{tricksbox}

\subsubsection*{5. Trace}

\begin{tracebox}[{Three reversals on Sample 1: $[1,2,3,4,5,6,7]$, $k = 3$}]
\centering
\begin{tikzpicture}[x=1cm, y=1cm]
    \dpindex{ix}{3.15}{index}{0}{6}
    \dparray{s0}{2.5}{start}{1/acell, 2/acell, 3/acell, 4/acell, 5/acell, 6/acell, 7/acell}
    \dparray{s1}{1.1}{rev $[0,7)$}{7/awrite, 6/awrite, 5/awrite, 4/acell, 3/awrite, 2/awrite, 1/awrite}
    \dparray{s2}{-0.3}{rev $[0,3)$}{5/awrite, 6/acell, 7/awrite, 4/acell, 3/acell, 2/acell, 1/acell}
    \dparray{s3}{-1.7}{rev $[3,7)$}{5/akeep, 6/akeep, 7/akeep, 1/akeep, 2/akeep, 3/akeep, 4/akeep}
    \foreach \x/\y in {0/6, 1/5, 2/4}{ \draw[dep, <->, thin] (s1\x.north) to[bend left=22] (s1\y.north); }
    \draw[dep, <->, thin] (s20.north) to[bend left=30] (s22.north);
    \foreach \x/\y in {3/6, 4/5}{ \draw[dep, <->, thin] (s3\x.north) to[bend left=30] (s3\y.north); }
    \node[note, anchor=west, text width=3.3cm, align=left] at (6.8,1.1) {3 swaps; the middle cell 4 stays};
    \node[note, anchor=west, text width=3.3cm, align=left] at (6.8,-0.3) {1 swap; cell 1 stays};
    \node[note, anchor=west, text width=3.3cm, align=left] at (6.8,-1.7) {2 swaps; result};
\end{tikzpicture}
\figcap{Blue cells were written by the reversal of that row; the arcs show the swapped pairs. Six swaps in total.}

\medskip
\small
\begin{tabular}{@{}llccl@{}}
\toprule
\textbf{call} & \textbf{range} & \textbf{pairs $(l, r)$ swapped} & \textbf{swaps} & \textbf{array after the call} \\
\midrule
$k \leftarrow 3 \bmod 7$ & --- & --- & 0 & $[1,2,3,4,5,6,7]$ \\
\texttt{rev(0, 7)} & $[0,7)$ & $(0,6), (1,5), (2,4)$ & 3 & $[7,6,5,4,3,2,1]$ \\
\texttt{rev(0, 3)} & $[0,3)$ & $(0,2)$ & 1 & $[5,6,7,4,3,2,1]$ \\
\texttt{rev(3, 7)} & $[3,7)$ & $(3,6), (4,5)$ & 2 & $[5,6,7,1,2,3,4]$ \\
\bottomrule
\end{tabular}
\end{tracebox}

\begin{tracebox}[{The cycle method on $[10,20,30,40,50,60]$, $k = 2$}]
\centering
\begin{tikzpicture}[x=1cm, y=1cm]
    \dparray{t0}{2.5}{start}{10/acell, 20/acell, 30/acell, 40/acell, 50/acell, 60/acell}
    \dparray{t1}{0.9}{cycle of 0}{50/awrite, 20/acell, 10/awrite, 40/acell, 30/awrite, 60/acell}
    \dparray{t2}{-0.4}{cycle of 1}{50/akeep, 60/akeep, 10/akeep, 20/akeep, 30/akeep, 40/akeep}
    \foreach \x/\y in {0/2, 2/4}{ \draw[dep, thin] (t0\x.north) to[bend left=35] (t0\y.north); }
    \draw[dep, thin] (t04.south) to[bend left=12] (t00.south);
    \node[note, anchor=west, text width=3.6cm, align=left] at (5.9,1.55) {$\gcd(6,2) = 2$: start at 0, then at 1};
    \node[note, anchor=west, text width=3.6cm, align=left] at (5.9,-0.4) {$= [50,60,10,20,30,40]$, the right rotation by 2};
\end{tikzpicture}
\figcap{After the first cycle the even positions are final (blue); the second cycle fixes the odd ones.}

\medskip
\small
\begin{tabular}{@{}cccccl@{}}
\toprule
\textbf{cycle} & \texttt{cur} $\to$ \texttt{nxt} & \textbf{dropped} & \textbf{picked up} & & \textbf{array} \\
\midrule
0 & (load) & --- & 10 & & $[10,20,30,40,50,60]$ \\
0 & $0 \to 2$ & 10 & 30 & & $[10,20,10,40,50,60]$ \\
0 & $2 \to 4$ & 30 & 50 & & $[10,20,10,40,30,60]$ \\
0 & $4 \to 0$ & 50 & (10, discarded) & & $[50,20,10,40,30,60]$ \\
1 & (load) & --- & 20 & & $[50,20,10,40,30,60]$ \\
1 & $1 \to 3$ & 20 & 40 & & $[50,20,10,20,30,60]$ \\
1 & $3 \to 5$ & 40 & 60 & & $[50,20,10,20,30,40]$ \\
1 & $5 \to 1$ & 60 & (20, discarded) & & $[50,60,10,20,30,40]$ \\
\bottomrule
\end{tabular}
\end{tracebox}

\subsubsection*{6. Test suite}

\begin{testbox}
\centering\small
\begin{tabular}{@{}llll@{}}
\toprule
\textbf{nums} & $k$ & \textbf{Why it is interesting} & \textbf{Expected} \\
\midrule
$[1,2,3,4,5,6,7]$ & 3 & official Sample 1 & $[5,6,7,1,2,3,4]$ \\
$[-1,-100,3,99]$ & 2 & official Sample 2, negative values & $[3,99,-1,-100]$ \\
$[42]$ & 100000 & $n = 1$, huge $k$: must not crash & $[42]$ \\
$[1,2,3]$ & 0 & no rotation & $[1,2,3]$ \\
$[1,2,3]$ & 3 & $k = n$ & $[1,2,3]$ \\
$[1,2,3]$ & 4 & $k = n+1$, same as $k = 1$ & $[3,1,2]$ \\
$[1,2,3,4,5,6]$ & 4 & $\gcd = 2$ (cycle method needs two starts) & $[3,4,5,6,1,2]$ \\
$n = 10^5$ & 99999 & timing: the brute force does $10^{10}$ moves & left rotation by 1 \\
\bottomrule
\end{tabular}

\medskip
All three programs above (brute force, reversals, cycles) agree on both official samples and on 1000 random tests with $1 \le n \le 9$, $0 \le k \le 25$, values in $[-5,5]$ (\texttt{tools/stress.py}).
\end{testbox}

\subsubsection*{7. Common mistakes}

\begin{warningbox}[Mistake 1 --- no $k \bmod n$]
\begin{lstlisting}[style=snippet]
rev(nums, 0, k);        // k = 100000, n = 1: indices far outside the array
k %= n;                 // RIGHT: first line of the function
\end{lstlisting}
\end{warningbox}

\begin{warningbox}[Mistake 2 --- the reversals in the left-rotation order]
``reverse $[0,k)$, reverse $[k,n)$, reverse all'' gives $[4,5,6,7,1,2,3]$ on Sample 1 --- the \emph{left} rotation by 3. If the output looks rotated the wrong way, swap the order or use $k \leftarrow n - k$.
\end{warningbox}

\begin{warningbox}[Mistake 3 --- assigning a local copy]
\begin{lstlisting}[style=snippet]
vector<int> b(n); /* fill b */  nums.swap(b);   // RIGHT: the caller sees the new contents
vector<int> b(n); /* fill b */  auto c = b;      // WRONG: nums itself was never changed
\end{lstlisting}
The function returns \texttt{void}; only changes made to \texttt{nums} (taken by reference) count.
\end{warningbox}

\begin{warningbox}[Mistake 4 --- a single cycle from index 0]
Following $0 \to k \to 2k \to \cdots$ only once rotates just the positions $\equiv 0 \pmod{\gcd(n,k)}$. With $n = 6$, $k = 2$ the odd positions stay untouched: $[50,20,10,40,30,60]$ instead of $[50,60,10,20,30,40]$.
\end{warningbox}

\subsubsection*{8. Takeaways}

\begin{itemize}[leftmargin=*]
    \item A rotation is a swap of two blocks, $AB \to BA$; three reversals do it in place with at most $n$ swaps.
    \item Reduce the shift modulo $n$ before anything else; test $k = 0$, $k = n$ and $n = 1$.
    \item Moving values in place needs an order that never overwrites an unread value: reversals swap pairs, cycles carry one spare value.
    \item The cycle structure of $i \to (i+k) \bmod n$ has $\gcd(n,k)$ cycles --- a fact that reappears whenever something jumps by $k$ around a circle.
\end{itemize}

\begin{variantsbox}
\variant{Rotate List}{\exmedium\ \lc{Medium}\ \varlink{Rotate List}{LeetCode 61}{https://leetcode.com/problems/rotate-list/}}
  {the sequence is a singly linked list; rotate it to the right by $k$ ($k \le 2\cdot10^9$).}
  {model: the same $AB \to BA$, but blocks are cut and re-linked instead of moved; $k \bmod$ length still first.}
  {algorithm: find the length and the tail, link the tail to the head, cut after node $n - (k \bmod n)$; $\bigO{n}$, samples $\to [4,5,1,2,3]$ and $[2,0,1]$.}
\variant{Shift 2D Grid}{\exeasy\ \lc{Easy}\ \varlink{Shift 2D Grid}{LeetCode 1260}{https://leetcode.com/problems/shift-2d-grid/}}
  {an $m \times n$ grid is shifted $k$ times; each shift moves every cell one place along the row-major order, wrapping around.}
  {model: the row-major reading of the grid is a 1D array of length $mn$ rotated right by $k$.}
  {algorithm: cell $i$ (row-major) goes to $(i + k) \bmod mn$, i.e.\ row $j / n$, column $j \bmod n$; all three official samples reproduced.}
\variant{Rotate left}{\varnone}
  {rotate to the left by $k$ (the first $k$ elements go to the end).}
  {model: left by $k$ = right by $n - (k \bmod n)$; position $i$ goes to $(i - k) \bmod n$.}
  {algorithm: reverse $[0,k)$, reverse $[k,n)$, reverse all; or \texttt{std::rotate(begin, begin + k, end)}.}
\variant{Rotate a subarray}{\varnone}
  {only the range $[l, r)$ is rotated right by $k$; the rest stays.}
  {model: $AB \to BA$ inside the range, with $|B| = k \bmod (r - l)$.}
  {algorithm: the same three reversals restricted to $[l, r)$; $\bigO{r - l}$ time, $\bigO{1}$ memory (E3.3 chooses the best such rotation).}
\end{variantsbox}
